\documentclass[aspectratio=169]{beamer}
\usepackage[utf8]{inputenc}
\usepackage[portuguese]{babel}
\usepackage{graphicx}
\usepackage{amsmath}
\usepackage{hyperref}

\usetheme{Madrid}
\usecolortheme{default}

\title[Anatomia do Google]{A Anatomia de um Motor de Busca Web de Larga Escala Hipertextual}
\subtitle{Google: Um Protótipo Revolucionário}
\author{Sergey Brin \and Lawrence Page}
\institute{Departamento de Ciência da Computação\\Stanford University}
\date{1998}

\begin{document}

\frame{\titlepage}

\begin{frame}{Sumário}
\tableofcontents
\end{frame}

\section{Introdução}

\begin{frame}{Contexto e Motivação}
\begin{itemize}
\item A Web cria novos desafios para recuperação de informação
\item Crescimento rápido: de 110.000 páginas (1994) para 100+ milhões (1997)
\item Motores de busca tradicionais retornam muitos resultados de baixa qualidade
\item Problema: usuários só conseguem avaliar os primeiros resultados
\end{itemize}

\vspace{0.5cm}
\begin{block}{Por que "Google"?}
Comum ortografia de \textit{googol} ($10^{100}$), refletindo o objetivo de construir motores de busca de escala muito grande
\end{block}
\end{frame}

\begin{frame}{Problemas dos Motores de Busca Existentes}
\begin{columns}
\column{0.5\textwidth}
\textbf{Desafios Técnicos:}
\begin{itemize}
\item Escalabilidade
\item Velocidade de indexação
\item Processamento eficiente
\item Milhões de consultas/dia
\end{itemize}

\column{0.5\textwidth}
\textbf{Desafios de Qualidade:}
\begin{itemize}
\item "Resultados lixo"
\item Manipulação por anunciantes
\item Falta de precisão
\item Necessidade de alta relevância
\end{itemize}
\end{columns}

\vspace{0.5cm}
\begin{alertblock}{Exemplo Real}
Em 1997, apenas 1 dos 4 principais motores comerciais encontrava a si mesmo ao buscar seu próprio nome!
\end{alertblock}
\end{frame}

\section{Objetivos de Design}

\begin{frame}{Objetivos Principais do Google}
\begin{enumerate}
\item \textbf{Qualidade de Busca Melhorada}
\begin{itemize}
\item Alta precisão (relevância dos resultados)
\item Uso de estrutura hipertextual
\item Exploração de links e texto âncora
\end{itemize}

\item \textbf{Pesquisa Acadêmica}
\begin{itemize}
\item Publicação aberta de detalhes técnicos
\item Ambiente para experimentos em larga escala
\item Acesso a dados de uso reais
\end{itemize}

\item \textbf{Escalabilidade}
\begin{itemize}
\item Sistema utilizável por números razoáveis de pessoas
\item Suporte a pesquisas inovadoras
\item Processamento eficiente de grandes volumes
\end{itemize}
\end{enumerate}
\end{frame}

\section{PageRank}

\begin{frame}{PageRank: Trazendo Ordem à Web}
\begin{block}{Conceito Fundamental}
Medida objetiva da importância de citação de uma página que corresponde bem com a ideia subjetiva de importância das pessoas
\end{block}

\vspace{0.3cm}
\textbf{Princípios:}
\begin{itemize}
\item Utiliza o grafo de citações (links) da Web
\item Não conta todos os links igualmente
\item Normaliza pelo número de links em uma página
\item Propagação recursiva de pesos através da estrutura de links
\end{itemize}

\vspace{0.3cm}
\begin{exampleblock}{Analogia}
Páginas bem citadas de muitos lugares na Web valem a pena olhar. Páginas com apenas uma citação do Yahoo! também valem a pena!
\end{exampleblock}
\end{frame}

\begin{frame}{Fórmula do PageRank}
\begin{block}{Definição Matemática}
Assumindo que a página $A$ tem páginas $T_1...T_n$ que apontam para ela:

$$PR(A) = (1-d) + d\left(\frac{PR(T_1)}{C(T_1)} + ... + \frac{PR(T_n)}{C(T_n)}\right)$$

onde:
\begin{itemize}
\item $d$ = fator de amortecimento (geralmente 0.85)
\item $C(A)$ = número de links saindo da página $A$
\item PageRanks formam uma distribuição de probabilidade
\end{itemize}
\end{block}

\vspace{0.3cm}
\textbf{Propriedades:}
\begin{itemize}
\item Calculável por algoritmo iterativo simples
\item Corresponde ao autovetor principal da matriz de links normalizada
\item 26 milhões de páginas processadas em poucas horas
\end{itemize}
\end{frame}

\begin{frame}{Justificativa Intuitiva do PageRank}
\begin{block}{Modelo do "Surfista Aleatório"}
\begin{itemize}
\item Usuário recebe uma página aleatória
\item Continua clicando em links
\item Nunca usa o botão "voltar"
\item Eventualmente se entedia e começa em outra página aleatória
\item \textbf{PageRank = probabilidade do surfista visitar a página}
\end{itemize}
\end{block}

\vspace{0.3cm}
\begin{columns}
\column{0.5\textwidth}
\textbf{Vantagens:}
\begin{itemize}
\item Difícil de manipular
\item Personalização possível
\item Reflexo do comportamento real
\end{itemize}

\column{0.5\textwidth}
\textbf{Fator de Amortecimento ($d$):}
\begin{itemize}
\item Probabilidade de continuar clicando
\item Pode ser aplicado seletivamente
\item Permite personalização
\end{itemize}
\end{columns}
\end{frame}

\section{Texto Âncora}

\begin{frame}{Texto Âncora (Anchor Text)}
\begin{block}{Inovação Importante}
Associar o texto dos links com a página para a qual o link aponta, não apenas com a página que contém o link
\end{block}

\vspace{0.3cm}
\textbf{Vantagens:}
\begin{enumerate}
\item \textbf{Descrições mais precisas} - âncoras frequentemente descrevem melhor as páginas do que as próprias páginas
\item \textbf{Busca de conteúdo não-textual} - permite buscar imagens, programas e bancos de dados
\item \textbf{Expansão de cobertura} - páginas não rastreadas podem ser retornadas
\end{enumerate}

\vspace{0.3cm}
\begin{exampleblock}{Dados do Sistema}
No rastreamento de 24 milhões de páginas: \textbf{259 milhões de âncoras} indexadas
\end{exampleblock}
\end{frame}

\section{Arquitetura do Sistema}

\begin{frame}{Arquitetura de Alto Nível}
\begin{center}
\textbf{Fluxo Principal de Dados:}
\end{center}

\begin{enumerate}
\item \textbf{URLserver} $\rightarrow$ envia listas de URLs para rastreadores
\item \textbf{Crawlers} (rastreadores) $\rightarrow$ baixam páginas Web
\item \textbf{StoreServer} $\rightarrow$ comprime e armazena páginas (repositório)
\item \textbf{Indexer} $\rightarrow$ processa documentos, cria hits e arquivo de âncoras
\item \textbf{URLResolver} $\rightarrow$ converte URLs e cria banco de links
\item \textbf{Sorter} $\rightarrow$ reorganiza dados por wordID (índice invertido)
\item \textbf{Searcher} $\rightarrow$ responde consultas usando índice e PageRank
\end{enumerate}

\vspace{0.3cm}
\begin{block}{Linguagens de Implementação}
C/C++ (eficiência) e Python (rastreadores) - executável em Solaris ou Linux
\end{block}
\end{frame}

\begin{frame}{Estruturas de Dados Principais}\footnotesize
\begin{columns}
\column{0.5\textwidth}
\textbf{Bigfiles:}
\begin{itemize}\itemsep=0cm
\item Arquivos virtuais
\item Múltiplos sistemas de arquivo
\item Suporte a compressão
\end{itemize}

\vspace{0.3cm}
\textbf{Repositório:}
\begin{itemize}\itemsep=0cm
\item HTML comprimido
\item ~50\% do armazenamento
\item Concatenação com cabeçalhos
\end{itemize}

\vspace{0.3cm}
\textbf{Índice de Documentos:}
\begin{itemize}\itemsep=0cm
\item Largura fixa (ISAM)
\item Ordenado por docID
\item Status, ponteiros, checksums
\end{itemize}

\column{0.5\textwidth}
\textbf{Hit Lists:}
\begin{itemize}\itemsep=0cm
\item Ocorrências de palavras
\item Posição, fonte, capitalização
\item Codificação compacta (2 bytes/hit)
\item Maior parte do espaço usado
\end{itemize}

\vspace{0.3cm}
\textbf{Índices:}
\begin{itemize}\itemsep=0cm
\item Forward: parcialmente ordenado
\item 64 barris por wordID
\item Inverted: processado pelo sorter
\item Doclist para cada palavra
\end{itemize}

\vspace{0.3cm}
\textbf{Otimização:}
\begin{itemize}\itemsep=0cm
\item Evitar disk seeks (10ms)
\item Estruturas eficientes
\item Processamento in-place
\end{itemize}
\end{columns}
\end{frame}

\begin{frame}{Sistema de Rastreamento}
\begin{block}{Desafios}
\begin{itemize}
\item Interação com centenas de milhares de servidores
\item Questões de desempenho e confiabilidade
\item Aspectos sociais e éticos
\end{itemize}
\end{block}

\vspace{0.3cm}
\textbf{Implementação:}
\begin{itemize}
\item Sistema distribuído em Python
\item Tipicamente 3 rastreadores paralelos
\item ~300 conexões abertas simultaneamente por rastreador
\item Cache DNS em cada rastreador
\item I/O assíncrono para gerenciar eventos
\end{itemize}
\end{frame}

\begin{frame}{Sistema de Rastreamento}

\vspace{0.3cm}
\begin{exampleblock}{Performance}
Pico: \textbf{100+ páginas/segundo} (4 rastreadores)\\
Últimos 11 milhões: \textbf{63 horas} = 4 milhões páginas/dia
\end{exampleblock}
\end{frame}

\section{Resultados}

\begin{frame}{Estatísticas do Sistema}
\begin{columns}
\column{0.5\textwidth}
\textbf{Armazenamento:}
\begin{itemize}
\item Páginas baixadas: 147.8 GB
\item Repositório comprimido: 53.5 GB
\item Índice invertido curto: 4.1 GB
\item Índice invertido completo: 37.2 GB
\item Léxico: 293 MB
\item Índice de documentos: 9.7 GB
\item Base de links: 3.9 GB
\item \textbf{Total: 55.2 GB (sem repo)}
\end{itemize}

\column{0.5\textwidth}
\textbf{Páginas Web:}
\begin{itemize}
\item Páginas baixadas: 24 milhões
\item URLs vistas: 76.5 milhões
\item Endereços email: 1.7 milhões
\item Erros 404: 1.6 milhões
\end{itemize}

\vspace{0.3cm}
\textbf{Performance:}
\begin{itemize}
\item Download: 9 dias (total)
\item Indexação: 54 páginas/seg
\item Ordenação: 24 horas (4 máquinas)
\item Consultas: 1-10 segundos
\end{itemize}
\end{columns}
\end{frame}

\begin{frame}{Exemplo de Resultados: "bill clinton"}
\begin{block}{Qualidade dos Resultados}
\textbf{Características demonstradas:}
\begin{itemize}
\item Agrupamento por servidor (whitehouse.gov)
\item Uso de texto âncora (primeiro resultado sem título)
\item Endereços de email retornados (5º resultado)
\item Todos os resultados de alta qualidade
\item PageRanks altos (percentagens mostradas)
\item Nenhum link quebrado
\item Proximidade de palavras respeitada
\end{itemize}
\end{block}

\vspace{0.3cm}
\begin{alertblock}{Comparação com Concorrentes}
Motores comerciais não retornavam resultados de whitehouse.gov, muito menos os corretos!
\end{alertblock}
\end{frame}

\begin{frame}{Comparação de Qualidade}
\begin{exampleblock}{Problema Real Observado}
Um motor de busca importante retornou "Bill Clinton Joke of the Day" como primeiro resultado para busca "Bill Clinton"
\end{exampleblock}

\vspace{0.3cm}
\textbf{Diferenciais do Google:}
\begin{enumerate}
\item \textbf{PageRank} - avaliação de qualidade das páginas
\item \textbf{Texto âncora} - descrição do que o link aponta
\item \textbf{Informação de proximidade} - relevância aumentada
\item \textbf{Informação de fonte} - tamanho e capitalização
\end{enumerate}

\vspace{0.3cm}
\begin{block}{Processo de Ranking}
Sistema de contagens $\rightarrow$ tabelas de lookup $\rightarrow$ rank final\\
Muitos parâmetros ajustáveis + sistema de feedback
\end{block}
\end{frame}

\section{Conclusões}

\begin{frame}{Contribuições Principais}
\begin{enumerate}
\item \textbf{PageRank} - algoritmo revolucionário de ranking baseado em links
\item \textbf{Uso extensivo de texto âncora} - melhora relevância e cobertura
\item \textbf{Arquitetura escalável} - design eficiente para crescimento da Web
\item \textbf{Sistema completo} - rastreamento, indexação e busca integrados
\item \textbf{Pesquisa acadêmica aberta} - publicação de detalhes técnicos
\end{enumerate}

\vspace{0.5cm}
\begin{block}{Resultados Alcançados}
\begin{itemize}
\item Protótipo com 24 milhões de páginas
\item Qualidade superior aos motores comerciais (avaliação subjetiva)
\item Sistema eficiente e escalável
\item Ferramenta de pesquisa para comunidade acadêmica
\end{itemize}
\end{block}
\end{frame}

\begin{frame}{Trabalho Futuro}
\begin{columns}
\column{0.5\textwidth}
\textbf{Melhorias de Eficiência:}
\begin{itemize}
\item Cache de consultas
\item Alocação inteligente de disco
\item Subíndices para termos comuns
\item Expansão para 100 milhões páginas
\end{itemize}

\vspace{0.3cm}
\textbf{Novas Funcionalidades:}
\begin{itemize}
\item Operadores booleanos
\item Negação e stemming
\item Feedback de relevância
\item Clustering avançado
\end{itemize}

\column{0.5\textwidth}
\textbf{Pesquisa Avançada:}
\begin{itemize}
\item Algoritmos de atualização inteligentes
\item Uso de proxy caches
\item Personalização de PageRank
\item Contexto do usuário
\item Sumarização de resultados
\end{itemize}

\vspace{0.3cm}
\textbf{Estrutura de Links:}
\begin{itemize}
\item Texto ao redor dos links
\item Extensões do PageRank
\item Novos usos da hipertextualidade
\end{itemize}
\end{columns}
\end{frame}

\begin{frame}{Visão Final}
\begin{block}{Ambiente Rico para Pesquisa}
Um motor de busca Web é um ambiente muito rico para ideias de pesquisa. O Google fornece infraestrutura para experimentação em larga escala.
\end{block}

\vspace{0.3cm}
\textbf{Objetivos de Longo Prazo:}
\begin{itemize}
\item Tornar-se um recurso para pesquisadores mundialmente
\item Estimular a próxima geração de tecnologia de busca
\item Manter foco em qualidade enquanto escala
\item Ambiente tipo "Spacelab" para experimentos em dados Web
\end{itemize}

\vspace{0.5cm}
\begin{center}
\Large{\textbf{Google: Trazendo ordem à Web através da estrutura hipertextual}}
\end{center}
\end{frame}

\begin{frame}{Agradecimentos e Suporte}
\textbf{Contribuidores Críticos:}
\begin{itemize}
\item Scott Hassan e Alan Steremberg
\end{itemize}

\vspace{0.3cm}
\textbf{Orientação e Discussões:}
\begin{itemize}
\item Hector Garcia-Molina, Rajeev Motwani
\item Jeff Ullman, Terry Winograd
\item Grupo WebBase
\end{itemize}

\vspace{0.3cm}
\textbf{Apoio Financeiro e Equipamentos:}
\begin{itemize}
\item Stanford Integrated Digital Library Project
\item National Science Foundation (IRI-9411306)
\item DARPA, NASA, Interval Research
\item Doadores: IBM, Intel, Sun
\end{itemize}
\end{frame}

\begin{frame}


\vspace{1cm}
\normalsize
\textbf{Disponível em:} http://google.stanford.edu

\vspace{0.5cm}
Sergey Brin e Lawrence Page\\
Computer Science Department\\
Stanford University
\end{frame}

\end{document}